\begin{problem}[Theorem 1. The Law of Sines:]
    Given a triangle with angle-side opposite pairs $(\alpha, a), (\beta, b), (\gamma, c)$, the following ratios hold:
    
    \[ \frac{sin(\alpha)}{a} = \frac{sin(\beta)}{b} = \frac{sin(\gamma)}{c}\]
    
    or, equivalently, 
    
    \[ \frac{a}{sin(\alpha)} = \frac{b}{sin(\beta)} = \frac{c}{sin(\gamma)}\]
\end{problem}

\begin{problem}[Theorem 2.]
    Suppose $(\alpha, a)$ and $(\gamma, c)$ are intended to be angle-side pairs in a triangle where $\alpha, a$ and $c$ are given. Let $h = c \cdot \sin(\alpha)$
    
    \begin{itemize}
        \item If $a < h$, then no triangle exists which satisfies the given criteria 
        \item If $a = h$, then $\gamma = \frac{\pi}{2}$ so exactly one (right) triangle exists which satisfies the criteria.
        \item If $h < a < c$, then two distinct triangles exist which satisfy the given criteria.
        \item if $a \geq c$, then $\gamma$ is acute and exactly one triangle exists which satisfies the given criteria.
    \end{itemize}
\end{problem}

\begin{problem}[Theorem 3. Law of Cosines:]
    Given a triangle with angle-side opposite pairs $(\alpha, a), (\beta, b), (\gamma, c)$, the following equations hold:
    
    \begin{align}
        \notag a^2 = b^2 + c^2 - 2 \cdot b \cdot c \cdot cos(\alpha) && b^2 = a^2 + c^2 - 2 \cdot a \cdot c \cdot cos(\beta) && c^2 = a^2 + b^2 - 2 \cdot a \cdot b \cdot cos(\gamma)
    \end{align}
    
    or, solving for the cosine in each equation, we have
    
    \begin{align}
        \notag \cos(\alpha) = \frac{b^2 + c^2 - a^2}{2bc} && \cos(\beta) = \frac{a^2 + c^2 - b^2}{2ac} && \cos(\gamma) = \frac{a^2 + b^2 - c^2}{2ab}
    \end{align}
\end{problem}

\begin{problem}[Converting Between Radians and Degrees:]
    Here for convenience.
    
    Converting from degrees to radians,
    
    \[1\degree \cdot \frac{1\pi}{180\degree} = \text{Radians}\]
    
    Converting from radians to degrees,
    
    \[1 \; \text{radian} \cdot \frac{180\degree}{\pi} = \text{Degrees}\]
\end{problem}